\section{Dividends: what the carry contains}
\label{sec:fwddivs}

So far the pair $(F,D)$ has been estimated without asking what makes
$F$ differ from spot in the first place.  This section opens that
box, because two later needs depend on it: the borrow of
\cref{sec:fwdborrow} is defined as what remains of the carry after
dividends are accounted for, and Chapter~7's early-exercise removal
takes a dividend model as its input.

Call \emph{carry} the growth rate the market charges between spot and
forward: the number $\log(F/S)/t$.  For a stock that pays nothing,
carry is just the rate --- $F=Se^{rt}$, because the alternative to
agreeing on delivery is to buy the stock today with borrowed money,
and the borrowing costs $r$.  Dividends change the arithmetic: whoever
holds the stock until expiry \emph{collects} them, and the forward
buyer does not, so their value is subtracted from what delivery is
worth.  How exactly it is subtracted depends on how the dividends are
modeled, and there are three standard conventions.

\emph{Discrete cash.}  The stock will pay fixed dollar amounts $d_i$
at \emph{ex-dates} $t_i\le t$ --- the dates from which a buyer of the
stock no longer receives the next dividend, so the price drops by
roughly the payment as each arrives.  (Both symbols are always
subscripted; the upright $\dd$ of an integral is unrelated.)  Replicate delivery: buy the stock
today for $S$, borrow the money, hold to expiry.  Along the way each
dividend $d_i$ arrives and can be lent out from its ex-date onward,
growing to $d_ie^{r(t-t_i)}$ by expiry.  The all-in cost of delivering
the stock at expiry is therefore the financed purchase minus the
collected, reinvested dividends:
\begin{equation}
F=Se^{rt}-\sum_i d_i\,e^{r(t-t_i)}
=\Big(S-\sum_i d_i\,e^{-rt_i}\Big)e^{rt}.
\label{eq:fwdescrow}
\end{equation}
The second form is the useful one: spot minus the \emph{present value}
of the dividend schedule, carried forward at the rate.  This is called
the escrowed-dividend forward --- as if the dividends were set aside
in escrow at time zero.

\emph{Discrete proportional.}  Each ex-date the price drops not by a
fixed dollar amount but by a fixed \emph{fraction} $f_i$ of whatever
the stock is worth that day.  The replication changes in one place:
reinvest each dividend into the stock itself rather than lending it.
Follow one ex-date to see the bookkeeping.  Just before it, a share
is worth some price $S^-$; the holder receives $f_iS^-$ in cash, and
the price drops to $(1-f_i)S^-$.  Spending the cash on shares at the
dropped price buys
\[
\frac{f_iS^-}{(1-f_i)S^-}=\frac{f_i}{1-f_i}
\quad\text{new shares per share held,}
\]
multiplying the holding by
$1+\tfrac{f_i}{1-f_i}=\tfrac{1}{1-f_i}$.  Each ex-date therefore
cancels one factor $(1-f_i)$: start with $\prod_i(1-f_i)$ of a share
today and exactly one share exists at expiry.  The cost today is
$S\prod_i(1-f_i)$, and
\begin{equation}
F=S\,\prod_i\big(1-f_i\big)\,e^{rt}.
\label{eq:fwdprop}
\end{equation}

\emph{Continuous yield.}  Smear the proportional drops into a
constant leak $q_d$ (subscripted --- bare $q$ is Chapter~2's quantile
density): the limit of \eqref{eq:fwdprop} with many small events is
\begin{equation}
F=S\,e^{(r-q_d)\,t}.
\label{eq:fwdyield}
\end{equation}
This is the cheapest convention, and the one the running board has
been using since \cref{sec:fwdline}.

Which convention should a modeler pick?  At first sight it barely
matters.  \Cref{fig:fwddivs}(a) shows the chapter's running schedule
--- \$\MacFwdDivCashDollars\ per quarter on a \$100 stock, first
ex-date \MacFwdDivFirstDays\ days out --- through the lens of the
\emph{implied dividend yield}: the constant $q_d$ that would produce
the same forward at each maturity, $q_d(T)=r-\log(F(T)/S)/T$.  The
cash schedule draws a sawtooth.  Each vertical jump is one ex-date
entering the horizon: just after the first one enters, a single
\$\MacFwdDivCashDollars\ payment thirty days away annualizes to
\MacFwdDivToothPct\% --- violent at the short end, where dividing by a
small $T$ amplifies everything (the same $1/t$ the rate suffered in
\cref{sec:fwdident}).  By a year the teeth have shrunk toward the flat
equivalent yield of \MacFwdDivQeqPct\% drawn dashed, and the two
conventions price long-dated forwards almost alike.  If the term
structure of $F$ were all the convention touched, the choice would be
cosmetic.

It is not, and panel~(b) shows the quantity that separates the
conventions: what happens to the forward when \emph{spot moves}.  Differentiate each convention
in $S$ at a fixed schedule.  For the yield and proportional forms,
\eqref{eq:fwdyield} and \eqref{eq:fwdprop}, $F$ is proportional to
$S$, so
\[
\frac{\partial\log F}{\partial\log S}=1 :
\]
the forward rides spot one for one.  For the cash form, write
$\mathrm{PV}=\sum_id_ie^{-rt_i}$ for the schedule's present value;
then $F=(S-\mathrm{PV})\,e^{rt}$ with $\mathrm{PV}$ \emph{fixed in
dollars}, and
\begin{equation}
\frac{\partial\log F}{\partial\log S}
=S\,\frac{\partial}{\partial S}\log\big(S-\mathrm{PV}\big)
=\frac{S}{S-\mathrm{PV}}\;>\;1 .
\label{eq:fwdelas}
\end{equation}
A cash schedule does not scale with the stock, so on a rally the
forward outruns spot: the dividends take a fixed slice, and the
slice matters relatively less as the pie grows.  Panel~(b) plots this
\emph{spot elasticity} against maturity.  The cash staircase climbs
--- each step one more ex-date inside the horizon --- reaching
$S/(S-\mathrm{PV})=\MacFwdDivElasTwoYr$ at two years, where the
escrowed value is $\mathrm{PV}=\$\MacFwdDivPvTwoYr$; the dotted line
is the closed form
\eqref{eq:fwdelas} and the solid one a direct bump measurement, in
agreement.  The proportional and yield conventions sit at exactly
one, at every maturity.

\begin{figure}[htbp]
\centering
\paperfig{\textwidth}{fig_fwd_divs.pdf}
\caption{Dividend conventions on the running schedule
(\$\MacFwdDivCashDollars\ quarterly, first ex-date
\MacFwdDivFirstDays\ days out).  (a)~Implied dividend yield of the
cash schedule against maturity: a sawtooth, each tooth one ex-date
entering the horizon, spiking to \MacFwdDivToothPct\% at the short
end; dashed, the flat equivalent yield \MacFwdDivQeqPct\%.
(b)~The spot elasticity \eqref{eq:fwdelas}: cash climbs to
\MacFwdDivElasTwoYr\ by two years (dotted: closed form; solid:
measured by bumping spot); proportional and yield sit at exactly
one.}
\label{fig:fwddivs}
\end{figure}

Why does an elasticity deserve this much attention in a chapter about
statics?  Because the smile is drawn against $F$.  When spot moves and
the surface is transported to the new spot --- the scenario logic
Chapter~9 develops --- every strike's log-moneyness
$k=\log(K/F)$ moves by \emph{minus} however much $\log F$ moves, and
\eqref{eq:fwdelas} says the dividend convention decides that
amount.  Two modelers with identical fits
today, one carrying cash dividends and one a yield, will disagree
about which strike is at the money \emph{after} a two-percent rally
--- a risk disagreement, not a pricing one.  The dividend model is
thus a convention in the full sense this book gives the word: not
inferable from today's chain (panel~(a) showed the term structures
nearly agree), consequential for what the desk does next
(panel~(b)), and therefore to be stated, not assumed silently.

The carry now has two named parts --- rate and dividends --- and the
chapter can ask its last estimation question: whether what is left
over, the borrow, can be measured at all.  One tool is still missing:
how a carry error shows up \emph{in the smile}.  That tool is the
next section, and it doubles as the payoff of the whole chapter.
